% Section 4: Method.
% Every number below carries a `% src:` comment naming the committed file it
% comes from. Paths are relative to the repository root. No number here
% comes from a run that is not committed.

\section{Method}\label{sec:method}

\subsection{Observations}

Every experiment reduces to one kind of observation: a fact recorded about
an \emph{entity} at some moment, and a later check of whether it still
matched reality an interval later. We call the interval the \emph{lag} and
record each observation as a triple
\[
  (\text{entity},\ \text{lag},\ \text{still correct?}).
\]
An entity is whatever was observed: a company, an airport, an aircraft, a
vessel, a software package, a drug label or an electricity balancing
authority. ``Still correct'' is defined per fact, with an explicit
tolerance where the value is numeric (Table~\ref{tab:datasets}).
% src: calibration/fidcal.py (module docstring: "an (entity, lag_seconds, still_correct) triple")

The survival curve $S(\ell)$ is the fraction of observations at lag $\ell$
that are still correct. The \emph{empirical half-life} is the first lag at
which $S$ falls to $0.5$, linearly interpolated between the lags measured;
it is undefined when $S$ stays above $0.5$ over the observed range.
% src: calibration/fidcal.py survival_curve(), empirical_half_life()

\subsection{Datasets}

Table~\ref{tab:datasets} lists the seven sources. All are public, all are
read with the Python standard library, and each script writes its results
to \texttt{calibration/results/}.
% src: calibration/README.md "Running"

\begin{table}[t]
\centering
\small
\caption{Facts, ground truth and sample sizes. ``Entities'' and
``observations'' are as reported in each result file.}
\label{tab:datasets}
\begin{tabular}{p{2.6cm}p{3.6cm}p{3.2cm}rr}
\toprule
Source & Fact & Still correct means & Entities & Observations \\
\midrule
Wikidata, snapshots 2004--2016, lags quarterly to 10 years
  & the CEO of company $X$ is $Y$
  & $Y$ has no end date and no successor has started
  & 283 & 91{,}799 \\
METAR, 10 US airports, 2024
  & the temperature at airport $A$ is $T$
  & within $\pm 2\,^\circ$F
  & 10 & 1{,}575{,}505 \\
  & airport $A$'s flight category is $C$
  & same category (VFR/MVFR/IFR/LIFR)
  & 10 & 1{,}575{,}714 \\
  & airport $A$ is IFR or worse
  & still IFR or LIFR
  & 10 & 105{,}909 \\
NOAA AIS, six days of 2023, cargo ships and tankers
  & vessel $V$ is at position $P$
  & next report within 1\,km (also 10\,km)
  & 6{,}146 & 2{,}945{,}357 \\
BTS on-time data, all US domestic flights of 2024
  & the aircraft's $k$-th leg of the day arrives on schedule
  & within 15\,min, not cancelled or diverted; lag = legs flown before it
  & 6{,}112 & 7{,}046{,}561 \\
PyPI, monthly snapshots Jan 2015 -- Oct 2024
  & the latest version of package $X$ is $Y$
  & $Y$ is still the highest final, unyanked release
  & 915 & 1{,}016{,}717 \\
DailyMed label versions, 2008 onward, lags to 5 years
  & the warnings of drug label $X$ say $Z$
  & warnings text unchanged, word for word
  & 173 & 113{,}410 \\
EIA hourly demand, 14 balancing authorities, 2024
  & hourly demand in region $R$ is $D$
  & demand at the later hour within $\pm 5\%$ of $D$
  & 14 & 2{,}433{,}786 \\
\bottomrule
\end{tabular}
\end{table}
% src (rows, in order):
%   calibration/results/wikidata_ceo.md; calibration/wikidata_ceo.py SNAPSHOT_YEARS, LAGS
%   calibration/results/weather_metar.md (three sections); calibration/weather_metar.py STATIONS, YEAR
%   calibration/results/ships_ais.md "AIS 1 km, all vessels" and "AIS 10 km, all vessels"; calibration/ships_ais.py DAYS
%   calibration/results/flights_bts.md "Flight legs"
%   calibration/results/pypi_latest.md "all snapshots"; calibration/pypi_latest.py SNAPSHOTS
%   calibration/results/dailymed_warnings.md "all snapshots"; calibration/dailymed_warnings.py SNAPSHOTS, LAG_DAYS
%   calibration/results/eia_demand.md "EIA hourly demand"; calibration/eia_demand.py YEAR
% NOTE: docs/FINDINGS.md and calibration/README.md say "6,592 cargo ships and
% tankers"; the committed result says 6,146 entities. This draft uses 6,146.
% NOTE: calibration/README.md says "977 most downloaded packages" for PyPI;
% the committed result says 915 entities. This draft uses 915.

The AIS sample is the 15th of every other month in 2023, each observed for
12 hours. Vessels are split by what they were doing when observed:
underway (speed $\geq 1$ knot) or moored or at anchor.
% src: calibration/ships_ais.py module docstring and DAYS; calibration/README.md "Caveats"
For drug labels, openFDA serves only each label's current version: of 50
revised labels looked up by their identifier, none returned a past version.
Version history therefore comes from DailyMed.
% src: calibration/results/openfda_labels.md
The DailyMed sample is 400 prescription labels, of which 375 had warnings
text; 173 had a DailyMed version early enough to be followed for five
years.
% src: calibration/results/dailymed_warnings.md (last line); calibration/README.md "Caveats"

\subsection{The model}

Per fact type, we fit
\begin{equation}\label{eq:decay}
  S(\ell) = b + (s_0 - b)\, 2^{-\ell/h},
\end{equation}
where $s_0$ is the fraction correct at lag~0, $h$ the half-life, and $b$ the
\emph{baseline}: the fraction still correct once the fact is fully stale.
With $b = 0$ this is plain exponential decay to zero. Where a fact is
periodic we also fit an \emph{echo} model, which adds back the fresh fact's
agreement at every multiple $nP$ of a period $P$, with strength $r$, width
$w$, and a fade with half-life $h_r$ over periods:
\begin{equation}\label{eq:echo}
  S(t) = b + (s_0 - b)\Big[\, 2^{-t/h} + \sum_{n \geq 1} r\, 2^{-nP/h_r}\, 2^{-|t - nP|/w} \Big].
\end{equation}
% src: calibration/fidcal.py exp_model(); docs/FINDINGS.md "Model and method", "Finding 5"
Section~\ref{sec:model} defines both as fid computes them; the study's echo
model and fid's agree to $10^{-12}$.
% src: docs/FINDINGS.md "What changed in fid" row 5

\subsection{Fitting}

$s_0$ is fixed at the curve's value at lag~0. The half-life is found by
weighted least squares on the survival curve, each lag weighted by its
number of observations, with a refined logarithmic grid search (four
rounds of 201 points). With a baseline, $b$ is searched over
$\{0, 0.01, \ldots, 0.99\} \times s_0$, refitting $h$ for each value. If no
decay fits better than every finite half-life, $h = \infty$. The method is
deliberately simple and dependency-free; it is not a maximum-likelihood
fit.
% src: calibration/fidcal.py fit_exponential(), _fit_half_life()

The \emph{shape error} is the largest absolute gap between the empirical
curve and the fitted model over all measured lags. We call a fit poor when
it exceeds $0.10$.
% src: calibration/fidcal.py shape_error(), summary_markdown() ("poor fit" if > 0.10)

The parameters reported as ``fitted'' in Section~\ref{sec:results} are fitted on all
the data, to describe the curve. Every held-out score uses a separate fit on
the training half only.
% src: calibration/fidcal.py analyze(): s0/h/b from train_curve for Brier; s0_all/h_all for the reported fit

\subsection{Held-out scoring}

Entities are split 50/50 by the first byte of the SHA-256 of their
identifier, so no entity is in both halves. Models are fitted on one half
and scored on the other with the Brier score, the mean squared error of
the predicted probability (lower is better). Each model is compared with
two references:
\begin{itemize}
  \item a \textbf{constant}: the training half's overall fraction correct,
    whatever the age;
  \item the \textbf{declared lifetime}: the hand-picked half-life for that
    information type in fid's example registry
    (\texttt{examples/half-lives.yaml}), used as a declared lifetime would be,
    with $s_0 = 1$ and no baseline.
\end{itemize}
We also report reliability: predictions are bucketed into tenths, and the
mean prediction in each bucket is compared with the observed rate.
% src: calibration/fidcal.py split_train(), brier(), reliability(), analyze() ("The guess as fid would use it: source fidelity 1")

The echo test (Section~\ref{sec:results-echo}) uses every hourly lag from
0 to 14 days and two splits: by station (6 airports train, 4 test) and by
time (January--June train, July--December test).
% src: calibration/results/weather_periodic.md (opening paragraph)

\subsection{Rules fixed before the run}

Two experiments were scored against a pass rule written before they were
run.
\begin{itemize}
  \item \textbf{Echo.} The echo model is adopted for a fact only if it
    beats decay-to-baseline on held-out Brier in both splits, its
    reliability is no worse in both splits, and its largest gap from the
    curve is under $0.10$.
    % src: calibration/results/weather_periodic.md (rule lines); docs/FINDINGS.md "Finding 5" ("with the pass rule fixed before the run")
  \item \textbf{Re-check scheduling} (Section~\ref{sec:scheduling}). Its
    rule is stated with its results.
    % src: calibration/results/decisions.md "The rule"; docs/ROADMAP.md "Decision benchmark, 20 October" ("The rule set before running failed")
\end{itemize}
The rules were recorded in the project's own repository before the runs,
not in an external registry.
% src: docs/ROADMAP.md; git history of calibration/decisions.py and calibration/weather_periodic.py

\subsection{What we do not have}

The plan for this paper calls for bootstrap 95\% intervals over entities
for every headline number, and a time-based split for every dataset. The
committed results have intervals only for the re-check benchmark, and a
time split only for the echo tests. Both are listed as limitations in
Section~\ref{sec:limitations} until they are run.
% src: paper plan "Gaps to close first"; calibration/results/decisions.md (only file with 95% ranges); calibration/results/weather_periodic.md, eia_demand.md (only time splits)
